\documentclass[xcolor=dvipsnames, USenglish]{beamer}  %notes=show to print them in the generated pdf
% Packages for a reasonable beamer session
\usepackage[T1]{fontenc}
%\usepackage[ansinew]{inputenc}
\usepackage[utf8]{inputenc}
\usepackage{textcomp}
\usepackage{lmodern}
\usepackage{csquotes}
\usepackage[english]{babel}
%\usepackage{babel}

\usepackage{graphicx}

\usepackage{amsmath}
\usepackage{amsfonts}
\usepackage{amssymb}
\usepackage{amsthm}
\usepackage{bm}

\usepackage{booktabs}
\usepackage{tabularx}

\usepackage{hyperref}

\usepackage{ellipsis}

% Additional packages
\usepackage{graphicx}
\usepackage{subfigure}
\usepackage{xcolor}
\usepackage{minted}

\usepackage[style=authoryear, backend=biber]{biblatex}
\setbeamertemplate{itemize/enumerate body begin}{\setlength{\leftmargini}{1.5em}}
\renewcommand*{\nameyeardelim}{\addcomma\addspace}
\addbibresource{\jobname.bib}
\renewcommand{\footnotesize}{\tiny}

\usepackage{tikz,pgf,calc}
%\usetikzlibrary{matrix, shapes, positioning, calc,
%  decorations.pathreplacing, shapes.geometric, arrows}
\usetikzlibrary{shapes.geometric, arrows, calc}
\tikzstyle{startstop} = [rectangle, rounded corners, minimum width=3cm, minimum height=1cm, text centered, text width=1.5cm, draw=black, fill=red!30]
\tikzstyle{io} = [trapezium, trapezium left angle=70, trapezium right angle=110, minimum width=3cm, minimum height=1cm, text centered, draw=black, fill=blue!30]
\tikzstyle{process} = [rectangle, minimum width=3cm, minimum height=1cm, text centered, draw=black, fill=orange!30]
\tikzstyle{decision} = [diamond, minimum width=3cm, minimum height=0.5cm, text centered, draw=black, fill=green!30]
\tikzstyle{arrow} = [thick,->,>=stealth]

%% References
\newlength\leftsidebar
\makeatletter
\setlength\leftsidebar{\beamer@leftsidebar}
\makeatother

\usepackage[absolute,overlay]{textpos}
\newenvironment{reference}[2]{%
  \begin{textblock*}{\textwidth}(\leftsidebar+#1,\paperheight-#2)
      \scriptsize\bgroup\color{red!50!black}}{\egroup\end{textblock*}}

% Path to graphics
\graphicspath{{../img/}}

% Sources
\usepackage{setspace}
\newcommand{\source}[1]{\begin{spacing}{0.5}{\fontsize{5}{6}\selectfont source: \itshape {#1}}\end{spacing}}
 % PACKAGES
% Collection of useful mathematical symbols and commands
% Calculus
\newcommand{\ud}{\mathrm{d}}
\newcommand{\pder}[2]{\frac{\partial{#1}}{\partial{#2}}}
\newcommand{\dpder}[2]{\frac{\partial^2{#1}}{\partial{#2^2}}}
\newcommand{\sderp}[3]{\frac{\partial^2{#1}}{\partial{#2}\partial{#3}}}
\newcommand{\tder}[2]{\frac{\ud{#1}}{\ud{#2}}}
\newcommand{\rot}[1]{\nabla \times {#1}}
\newcommand{\diver}[1]{\nabla \cdot {#1}}
\newcommand{\definter}[4]{\int_{#1}^{#2} {#3}\ud {#4}}
\newcommand{\inter}[2]{\int {#1}\ud {#2}}
\newcommand{\braket}[2]{\langle {#1} , {#2} \rangle}
% Misc
\newcommand{\eval}[1]{\Big |_{#1}}
\newcommand{\bset}[1]{\big\lbrace {#1} \big\rbrace}
\newcommand{\stimes}[2]{{#1}\!\times\!{#2}}
\newcommand{\trp}{\top}
\newcommand{\preup}[2]{{}^{#2}\!{#1}}

% Operators
\DeclareMathOperator*{\armin}{arg\,min}
\DeclareMathOperator*{\armax}{arg\,max}
\DeclareMathOperator*{\rank}{rank}
\DeclareMathOperator*{\cov}{cov}
\DeclareMathOperator*{\nullsp}{null}
% Logicals
\newcommand{\suchthat}{\big \backslash \;}
   % SYMBOLS

% ----------- extra packages
\usepackage{../beamer_themes/beamerthemeEawag_blue} % Eawag style

% ----------- Extra symbols
\newcommand{\ccov}[1]{{\color{red}k}\left(#1\right)}
\newcommand{\cmean}[1]{{\color{blue}m}\left(#1\right)}
\newcommand{\sm}{\scalebox{0.5}{-1}}

% ----------- For boxed equations
\usepackage{amsmath}
\usepackage{empheq}
\usepackage[most]{tcolorbox}
\newtcbox{\mymath}[1][]{%
    nobeforeafter, math upper, tcbox raise base,
    enhanced, colframe=blue!30!black,
    colback=blue!30, boxrule=1pt,
    #1}



%----------------
% title information
\title{Master Thesis - Final Presentation}
\subtitle{Applications of Emulation to the Hydrology Field\\\textcolor{Red}{not really just hydrology, e.g. weir equation}}
\author[\texttt{sebastiano.rusca@eawag.ch}]{Sebastiano Rusca}
\institute[Eawag]{Eawag: Swiss Federal Institute of Aquatic Science
  and Technology}
\date[12.03.2018]{March 12, 2018}

% ====================================================================

\begin{document}

% ----------------
% Title frame
% load background for title
\setbeamertemplate{background}{
  \includegraphics[width=\paperwidth,height=\paperheight]
  {../beamer_themes/background_title_blue.png}}
{ \setbeamertemplate{footline}{} % no footer on title
  \begin{frame}
    \titlepage
  \end{frame}
}
% load background for all other slides
\setbeamertemplate{background}{
\includegraphics[width=\paperwidth,height=\paperheight]
{../beamer_themes/background_slides_blue.png}}
\setbeamertemplate{footline}[Sebastiano Rusca] % set footer
\addtocounter{framenumber}{-1}  % don't count title page

%%%%%%%%%%%%%%%%%%%%%%%%%%%%%%%%%%%%%%%%%%%%%%%
%----------------
\section{Introduction}

  \begin{frame}
    \frametitle{Outline}
    \begin{itemize}
      \item Master thesis goals
      \item What is emulation?
      \item Case studies
      \begin{itemize}
      %\itemsep0em
        \item A mechanistic emulator: \emph{learning the weir equation}
        \item A hyrological emulator: \emph{estimating the time-to-threshold}
      \end{itemize}
      \item Conclusions
      \item Outlook
     \end{itemize}
  \end{frame}

%%%%%%%%%%%%%%%%%%%%%%%%%%%%%%%%%%%%%%%%%%%%%%%%%%%%%%%%%%%%%%%%%%%%%%%%%%%%%%%%
% MASTER THESIS GOALS
  \begin{frame}
    \frametitle{Master thesis goals}
    \begin{itemize}
      \item Get an overview of emulation
      \begin{itemize}
        \item Learn about \emph{interpolation}, \emph{intrapolation} and \emph{extrapolation}
        \item Learn about \emph{Gaussian Processes}
      \end{itemize}
      \item Find an appropriate \emph{open source} hydrological simulator
      \begin{itemize}
      %\itemsep0em
        \item Learn its \emph{functioning} and \emph{customize} it for the required scope
        \item Produce \emph{test datasets} usable for emulation
      \end{itemize}
      \item Perform emulation with a \emph{didactic} example
      \begin{itemize}
      %\itemsep0em
        \item Mechanistic emulator: \emph{the weir equation}
      \end{itemize}
      \item Perform emulation in an \emph{hydrological} context
      \begin{itemize}
      %\itemsep0em
        \item Hydrological emulator: \emph{time-to-threshold}
      \end{itemize}
      \item Draw conclusions about emulation potential, especially applied to the \emph{hydrology},
            \emph{urban water} and \emph{water management} fields
    \end{itemize}
  \end{frame}

%%%%%%%%%%%%%%%%%%%%%%%%%%%%%%%%%%%%%%%%%%%%%%%%%%%%%%%%%%%%%%%%%%%%%%%%%%%%%%%%
% WHAT IS EMULATION?


%%%%%%%%%%%%%%%%%%%%%%%%%%%%%%%%%%%%%%%%%%%%%%%%%%%%%%%%%%%%%%%%%%%%%%%%%%%%%%%%
% CASE STUDY 1: OVERVIEW
\section{Case study 1: weir equation}

  \begin{frame}
    \frametitle{Mechanistic emulator: \emph{the weir equation}}
    \textbf{Goal}: learn the weir equation from simulated data\\
    \begin{alertblock}{Weir equation}
      \setlength\abovedisplayskip{0pt}
      \begin{equation*}
        Q = \frac{2}{3}\: \textcolor{red}{\mu}\: B_w\: \sqrt{2g}\: h_{\"u}^{\textcolor{red}{a}},
        \quad usu.\: \textcolor{red}{a = 3/2}
      \end{equation*}
    \end{alertblock}
  \end{frame}
% * The weir equation was partially derived empirically
% * It can be derived analytically by assuming no correction of velocity distribution
%   is needed.
% * In order to obtain actual discharge it has to be accounted for correction of
%   velocity distribution
% * This is done by adding a coefficient C_d (or \mu) in front of the equation, which
%   is dependent on the type and shape of weir
% * To obtain the value of this coefficient the equation has to be fitted to 
%   experimental data.


%%%%%%%%%%%%%%%%%%%%%%%%%%%%%%%%%%%%%%%%%%%%%%%%%%%%%%%%%%%%%%%%%%%%%%%%%%%%%%%%
% CASE STUDY 1: GRID CONVERGENCE STUDY


%%%%%%%%%%%%%%%%%%%%%%%%%%%%%%%%%%%%%%%%%%%%%%%%%%%%%%%%%%%%%%%%%%%%%%%%%%%%%%%%
% CASE STUDY 1: METHODOLOGY

%%%%%%%%%%%%%%%%%%%%%%%%%%%%%%%%%%%%%%%%%%%%%%%%%%%%%%%%%%%%%%%%%%%%%%%%%%%%%%%%
% CASE STUDY 1: RESULTS AND DISCUSSION

%%%%%%%%%%%%%%%%%%%%%%%%%%%%%%%%%%%%%%%%%%%%%%%%%%%%%%%%%%%%%%%%%%%%%%%%%%%%%%%%
% CASE STUDY 2: OVERVIEW
\section{Case study 2: time-to-threshold}

\begin{frame}
  \frametitle{Experiment overview}
  frame text
\end{frame}



%%%%%%%%%%%%%%%%%%%%%%%%%%%%%%%%%%%%%%%%%%%%%%%%%%%%%%%%%%%%%%%%%%%%%%%%%%%%%%%%
% THE END
  {
  \setbeamertemplate{background}{}
  \setbeamercolor{background canvas}{bg=black}
  \setbeamercolor{normal text}{fg=white}
  \usebeamercolor[fg]{normal text}
  \begin{frame}[plain]
    \centering
    \Large{\textbf{THE END}}\\
  \end{frame}
  }

%%%%%%%%%%%%%%%%%%%%%%%%%%%%%%%%%%%%%%%%%%%%%%%%%%%%%%%%%%%%%%%%%%%%%%%%%%%%%%%%
% ADDITIONAL SLIDES
% Links to repositories
\section{Backup slides}
  \begin{frame}
    \frametitle{Links to my repositories}
    \small{\url{https://bitbucket.org/binello7/fswof2d}}\\
    \small{\url{https://bitbucket.org/binello7/master_thesis}}\\
    \small{\url{https://bitbucket.org/binello7/master_thesis/wiki/Home}}
  \end{frame}

\end{document}

